\section{Introduction}
\label{sec:intro}

Agents and assistants increasingly act on records that change: a user's clearance is revoked, a ticket is
rerouted, a preference is replaced. The record usually keeps the history, and the model is told that only
the latest entry counts. Retrieval benchmarks ask whether the model can then report the current value, and
recent work shows that old values do intrude on such reports, especially after many overwrites
\citep{wang2026unable,guo2026context,chattaraj2026dual}. A deployed system, however, rarely asks the model
to restate a value. It asks the model to \emph{use} it. Two questions follow that retrieval accuracy does
not answer. When a model can state the current value correctly, does the obsolete value still change the
decision it takes? And if it does, is the obsolete value acting as a superseded assignment of the queried
attribute, or simply as a word that occurs in the context near the entity?

The second question matters because the remedies differ. If the interference is tied to supersession, it
is a failure of state tracking, and update-aware memory, historical marking or attention routing toward the
current entry \citep{guo2026context} are the natural fixes. If any nearby mention of a decision-relevant
word does the same, the problem is contextual entrainment \citep{niu2025entrainment} reaching decisions.
Removing or marking stale entries would then leave much of it in place, and evaluations that only vary
update histories would misattribute it. The two accounts are not exclusive, and our design estimates
both.

\paragraph{Design.} We hold everything that should determine the decision fixed: the current state, the
policy that maps states to actions, the query, and the correct action. We edit one obsolete state. In one
member of each pair, the obsolete state implies the wrong action; in the other, it implies the correct
action. The difference in the model's preference for the wrong action between the two members is the
decision effect of the obsolete content, and the paired generated answers give its behavioral
counterpart (Figure~\ref{fig:example}). The same edit is repeated in matched constructions that keep the
word and its position but change its role: an assignment to a different attribute that was later updated,
a word the entity only mentioned, and a word attached to no entity. Before testing decisions, we check on
the same records that the model reports the current and the initial state correctly when asked.

\paragraph{Findings.} The experiment ran under a protocol frozen before the confirmation data were
generated, on a held-out prompt template, with Qwen3-8B. A larger model from the same family,
Qwen3-14B, was run on identical rows under a dated amendment, and Gemma 3 4B was attempted. The routing
task meets every competence criterion on the confirmation data in both Qwen models; the access task does not, and its estimates
are reported second and labeled. Three conclusions hold for both confirmed models, with one
model-specific difference:
\begin{enumerate}
  \item \textbf{Obsolete states change current decisions in models that track them correctly.} In the
  routing histories where the model retrieved the current and initial states and decided correctly without
  history, making the obsolete state imply the wrong action gave a net switch rate toward the wrong action
  of 0.158 (Qwen3-8B) and 0.165 (Qwen3-14B), almost never in reverse (\S\ref{sec:results-h1},
  \S\ref{sec:results-replication}).
  \item \textbf{Supersession adds decision pressure, and an incidental mention adds pressure of the same
  order.} In routing, a superseded assignment moved preference more than an updated assignment of an
  unrelated attribute, by 1.224 and 0.866 nats. The same word mentioned in passing by the same entity moved
  decisions as much as the superseded assignment in Qwen3-14B, and more in Qwen3-8B. Whether a mention
  \emph{exceeds} a superseded assignment therefore depends on the model, and in Qwen3-8B also on wording
  (\S\ref{sec:results-h2}).
  \item \textbf{A score-level sensitivity measure does not identify vulnerable items.} The edit's effect on
  easy retrieval queries was large, but in both models it added nothing beyond the pre-decision features
  under the specified logistic model (\S\ref{sec:results-h3}).
\end{enumerate}

\paragraph{What is and is not new.} That stale context interferes with answers is established
\citep{wang2026unable,guo2026context,qian2026forget,chao2026stale}, as is the fact that mentioned tokens
gain probability \citep{niu2025entrainment}. Our contribution is a controlled test that connects the two at
the level of decisions. We measure interference on policy decisions in open models, on items whose
current and historical states the model retrieves correctly. We use minimal-pair edits that hold the
correct action fixed, report both scores and generated answers, and include the matched controls needed
to ask whether supersession is the operative factor. We also report the negative result for a
sensitivity-based predictor and, in Appendix~\ref{app:history}, two measurement problems we found along
the way: chat-format events that the model never generates, and log-probabilities that change when tokens
are appended under 8-bit and bfloat16 inference.

\begin{figure*}[t]
\centering
\begin{prompt}
Read the dated record. Only the final assignment governs the decision.
Initially, EvbUnit0096A had routing state pending.        <- edited: pending | suspended
Initially, EvbUnit0096B had routing state suspended.
Later, EvbUnit0096A had archive tag restricted.
Finally, EvbUnit0096A had routing state restricted.
Later, EvbUnit0096B had archive tag restricted.
Finally, EvbUnit0096B had routing state pending.
Policy: state to bin: restricted -> B0, suspended -> B1, pending -> B2,
  active -> B3; bin to action: B0 -> EAST, B1 -> EAST, B2 -> WEST, B3 -> WEST
Apply the supplied rule to the final assignment of EvbUnit0096A; give its action.
Respond with exactly one label (EAST, WEST), with no other text.
\end{prompt}
\vspace{-0.5em}
{\small
\begin{tabular}{@{}ll@{}}
\toprule
Construction & Edited line (same word, same position) \\
\midrule
Superseded assignment & \lit{Initially, EvbUnit0096A had routing state pending.} \\
Other attribute, updated & \lit{Initially, EvbUnit0096A had archive tag pending.} \\
Mentioned by entity & \lit{Initially, EvbUnit0096A mentioned the word pending.} \\
Unattached word & \lit{Initially, the unassigned word was pending.} \\
\bottomrule
\end{tabular}}
\caption{One confirmation item (routing task, sequential order, held-out template; the annotation arrow is
not part of the prompt). The current state \lit{restricted} maps to \lit{EAST}. The obsolete state is
edited between \lit{pending} (implies \lit{WEST}, wrong) and \lit{suspended} (implies \lit{EAST}, correct);
everything else is identical. The table shows how the edited line reads in each construction. The decision
effect $D_c$ compares the model's preference for the wrong action across the two edits.}
\label{fig:example}
\end{figure*}
